%%%%%%%%%%%%%%%%%%%%%%%%%%%%%%%%%%%%%%%%%
% Medium Length Professional CV
% LaTeX Template
% Version 2.0 (8/5/13)
%
% This template has been downloaded from:
% http://www.LaTeXTemplates.com
%
% Original author:
% Trey Hunner (http://www.treyhunner.com/)
%
% Important note:
% This template requires the resume.cls file to be in the same directory as the
% .tex file. The resume.cls file provides the resume style used for structuring the
% document.
%
%%%%%%%%%%%%%%%%%%%%%%%%%%%%%%%%%%%%%%%%%

%----------------------------------------------------------------------------------------
%	PACKAGES AND OTHER DOCUMENT CONFIGURATIONS 
%----------------------------------------------------------------------------------------
 
\documentclass{resume} % Use the custom resume.cls style 

\usepackage[left=0.4 in,top=0.3 in,right=0.4 in,bottom=0.3in]{geometry} % Document margins

% \usepackage{biblatex}
% \addbibresource{publications.bib}
% \defbibheading{bibliography}[\refname]{}

\usepackage[sort]{natbib}
\bibliographystyle{unsrtnat}
\renewcommand{\bibsection}{}

\newcommand{\tab}[1]{\hspace{.2667\textwidth}\rlap{#1}}
\newcommand{\itab}[1]{\hspace{0em}\rlap{#1}}
\newcommand{\specialcell}[2][c]{%
  $\left.\begin{tabular}[#1]{@{}#1@{}}#2\end{tabular}\right\}$}
  \newcommand{\specialcellbracketless}[2][c]{%
  $\begin{tabular}[#1]{@{}#1@{}}#2\end{tabular}$}
\name{Noah Amsel} % Your name 
% \address{C-101, Dreamland Apartment, Backbone Park, B/H Balaji Hall, Rajkot-360004, Gujarat} % Your address 
%\address{123 Pleasant Lane \\ City, State 12345} % Your secondary addess (optional) 
\address{noah.amsel@nyu.edu \\ https://NoahAmsel.github.io/} % Your phone number and email

\begin{document}

\centerline{\textit{Last updated \today}}

%----------------------------------------------------------------------------------------
%	EDUCATION SECTION
%----------------------------------------------------------------------------------------

\begin{rSection}{Education}

{\textbf{Courant Institute, New York University}} \hfill {August 2022 - Present}
\\ 
Ph.D. in Computer Science \hfill {New York, NY}
\\
Advised by Joan Bruna and Chris Musco.

{\textbf{Yale University}} \hfill {August 2016 - May 2020}
\\ 
B.S. in Computer Science \& Mathematics, with distinction. \hfill {New Haven, CT}
\\
\textit{Magna cum laude}. Thesis: ``Online Vector Balancing in Practice,'' advised by Dan Spielman.

\end{rSection} 

%----------------------------------------------------------------------------------------
%	TECHNICAL STRENGTHS SECTION
%----------------------------------------------------------------------------------------

% \begin{rSection}{skills and INTERESTS}

% \begin{tabular}{ @{} >{\bfseries}l @{\hspace{6ex}} l }  
% Interests & Product Development, Design, Automobile, CAD/CAE, Finite Element Analysis, \\& Optimization, Fluid Mechanics, Robotics, Modeling and Simulation\\    
% Design Software & Basic AUTOCAD, CATIA V5, ANSYS (Static Structural, Transient Structural, \\& Static Thermal, Transient Thermal, Harmonic Response, Model analysis, Acoustic, Fluent), \\& OptimumLap,  MATLAB\\      
 
% \end{tabular}   

% \end{rSection}

%-------------------------------------------------------------------------------
%	PROJECTS

\begin{rSection}{Experience}

\begin{rSubsection}{Polymathic AI} {June 2024 - August 2024}{Summer Research Intern}{New York, NY}

\item Contributed to deep learning research on large foundation models for scientific data from heterogeneous domains.
\item Led project studying transfer capabilities of models for fluid dynamics problems described by differential equations.

\end{rSubsection}  
    
%------------------------------------------------    

\begin{rSubsection}{Qualcomm Technologies, Inc.} {November 2021 - July 2022}{Engineer, Corporate Research \& Development}{New York, NY}

\item Stayed on after Reservoir Labs was acquired by Qualcomm and turned into a new R\&D division.
\item Worked on software package for global non-convex optimization.
 
\end{rSubsection}  

%------------------------------------------------

\begin{rSubsection}{Reservoir Labs} {June 2020 - November 2021}{Research Engineer}{New York, NY}

\item Created algorithm for simulating network performance that achieved order-of-magnitude speed up.
\item Deployed our network modeling tool to DOE’s Energy Sciences Network.
\item Developed framework for designing provably efficient data center networks. Published it at ACM SIGCOMM.
 
\end{rSubsection}  

%------------------------------------------------

\begin{rSubsection}{Weizmann Institute of Science} {June - August 2019}{Kupcinet-Getz International Summer School, Research Fellow}{Rehovot, Israel}

\item Developed novel spectral method for fitting latent tree graphical models of DNA data.
\item Proved sufficient condition for the algorithm to succeed and finite sample guarantee.
\item Wrote open source implementation that scales to to very large problems with high accuracy.

\end{rSubsection}  

%------------------------------------------------

\begin{rSubsection}{Facebook, Inc.} {May - August 2018}{Software Engineering Intern}{New York, NY}

\item Built Bandwidth Estimation model for Adaptive Bitrate Streaming to improve mobile video quality.
\item Implemented and tested the model in C++ and Java; refactored ABR code.
\item Validated new model in production; found it reduced error by more than half (RMSE).
\item Received return offer.

\end{rSubsection}  

%------------------------------------------------

\begin{rSubsection}{Off the Hook, LLC} {May - July 2017}{Data Analyst / Software Developer}{New York, NY}

\item Developed computer gambling software in Python based on statistical analysis of horse racing data.
\item Evaluated wager opportunities using machine learning trained on data from hundreds of thousands of past races. 

\end{rSubsection}  

%------------------------------------------------

% \begin{rSubsection}{Yale Brain Function Laboratory} {Summer 2014, 2015}{Research Assistant}{New Haven, CT}

% \item Wrote MATLAB programs to analyze functional Near Infrared Spectroscopy data at the gigabyte scale.
% \item Introduced permutation testing for statistical significance; learned signal processing methods.
% \item Designed and conducted independent, award-winning research project that used human subjects to study the neural circuitry of interpersonal communication (Intel STS Semifinalist, 2016).

% \end{rSubsection}  

%------------------------------------------------

\end{rSection} 

%--------------------------------------------------------------------------------------
%   Research Publications 
%--------------------------------------------------------------------------------------
\begin{rSection}{ Publications } \itemsep -3pt        

\nocite{*}
\bibliography{publications.bib}

\end{rSection}

%----------------------------------------------------------------------------------------
%	AWARDS SECTION
%----------------------------------------------------------------------------------------

\begin{rSection}{Awards}

\begin{tabular}{ll}
    NSF Graduate Student Research Fellowship & 2024 \\
    Phi Beta Kappa & 2020 \\
    Intel Science Talent Search Semifinalist & 2016 
\end{tabular}

\end{rSection}

%----------------------------------------------------------------------------------------
%	TEACHING EXPERIENCE SECTION
%----------------------------------------------------------------------------------------

\begin{rSection}{Teaching Experience}

\begin{tabular}{l>{\em}ll}
    Algorithms (CPSC 365), Undergraduate Learning Assistant & Yale University & 2019 \\
    Volunteer Math Tutor & Top Honors & 2013--2016
\end{tabular}

\end{rSection}

%----------------------------------------------------------------------------------------
%	TALKS
%----------------------------------------------------------------------------------------

\begin{rSection}{Conference Talks}

\begin{tabular}{lll}
    NYAS ML Symposium & The Low-Rank Bottleneck in Attention & 2024 \\
    \specialcell[l]{SIAM Linear Algebra \\ Temple Num. Anal. Day} & Fixed-Sparsity Matrix Approximation from Matrix-Vector Products & 2024 \\
    SIAM-NNP & Near-Optimal Approximation of Matrix Functions by the Lanczos Method & 2023 \\
    ACM SIGCOMM & Designing Data Center Networks Using Bottleneck Structures & 2021 \\
    IEEE/ACM INDIS at SC20 & Computing Bottleneck Structures at Scale & 2020
\end{tabular}

\end{rSection}

%-------- TODO! remove me --------
\pagebreak
%-------- TODO! remove me --------
\begin{rSection}{Poster Presentations}

    \begin{tabular}{lll}
        \specialcell[l]{NeurIPS \\ ICERM \\ASE60 (MIT)\\Mihalis Fest (Columbia)} & Lanczos-FA is Nearly Krylov-Optimal for Rational Functions & \specialcellbracketless[l]{2024 \\ 2023 \\ 2023 \\ 2023} \\
        ACL & Finding Hierarchical Structure in Neural Stacks Using Unsupervised Parsing & 2019
    \end{tabular}
    
    \end{rSection}

\end{document}
